% 03-copy-task.tex ;  The copy task and results
%
% Scope: definition, why it's harder than adding, results with clear gradient decay

\section{Bài toán sao chép}
\label{sec:ch02-copy}

\index{copy task|(}%

\subsection{Định nghĩa}

Bài toán sao chép kiểm tra trí nhớ theo cách trực tiếp hơn: mạng phải đọc một
chuỗi, im lặng qua một khoảng trễ, rồi viết lại chính xác chuỗi đó.

Cụ thể, với tham số $T_{\text{mem}}$:
\begin{enumerate}
  \item $T_{\text{mem}}$ bước đầu tiên: mỗi bước là một vector 8-bit (mỗi
    bit là 0 hoặc 1, dạng one-hot; tức chỉ một bit bằng 1, bảy bit còn lại
    bằng 0).
  \item Một bước \textbf{delimiter}: tất cả bit dữ liệu bằng 0, kênh đánh
    dấu bằng 1.
  \item $T_{\text{mem}}$ bước tiếp theo: toàn bộ đầu vào bằng 0 (cả dữ liệu
    lẫn kênh đánh dấu). Mạng phải xuất ra đúng $T_{\text{mem}}$ vector 8-bit
    mà nó đã đọc ở bước~1.
\end{enumerate}

Tổng chiều dài chuỗi: $2T_{\text{mem}} + 1$. Mạng có $T_{\text{mem}}$ bước
để đọc, 1 bước để nhận tín hiệu ``bắt đầu chép'', và $T_{\text{mem}}$ bước
để viết lại; trong khi đầu vào là toàn số 0.

\begin{example}
Với $T_{\text{mem}} = 3$, dữ liệu là $[1,0,0,0,0,0,0,0]$,
$[0,0,1,0,0,0,0,0]$, $[0,0,0,0,1,0,0,0]$ (tương ứng với các ký hiệu A, C,
E trong bảng chữ cái 8 ký tự). Chuỗi đầu vào:
\[
\texttt{A C E | 0 0 0}
\]
với \texttt{|} là delimiter. Target: \texttt{A C E}.
\end{example}

Bài toán này khó hơn bài toán cộng vì mạng phải nhớ \textbf{toàn bộ} chuỗi
đầu vào (không chỉ hai giá trị), và phải duy trì trí nhớ đó qua $T_{\text{mem}}$
bước hoàn toàn im lặng.

\subsection{Kết quả}

Tôi huấn luyện RNN với $d_h = 32$, learning rate $0.01$, 2000 epoch, khởi tạo
$\mathcal{N}(0, 0.1)$. Hàm mất mát: cross-entropy trên từng bit đầu ra, chỉ
tính ở $T_{\text{mem}}$ bước cuối. Loss của dự đoán ngẫu nhiên trên 8 lớp là
$-\ln(1/8) \approx 2.08$ mỗi bước.

\begin{measured}{Copy task, $d_h = 32$, 2000 epoch}
\small
\begin{tabular}{@{}lrrrr@{}}
\toprule
$T_{\text{mem}}$ & Loss trước & Loss sau &
$\|\partial L / \partial \mat{W}_{hh}\|$ tại $t=1$ & tại $t=2T_{\text{mem}}+1$ \\
\midrule
5  & 8.74 & 8.34 & 0.233 & 0.033 \\
10 & 7.49 & 8.34 & 0.416 & 0.017 \\
20 & 7.91 & 8.37 & 0.832 & 0.033 \\
\bottomrule
\end{tabular}
\normalsize
\end{measured}

Ba điều đáng chú ý trong bảng này.

\textbf{Thứ nhất: gradient biến mất.} Tỉ lệ $\|\partial L / \partial
\mat{W}_{hh}\|$ giữa bước đầu tiên và bước cuối cùng tăng từ 7 lần ở
$T_{\text{mem}} = 5$ lên 25 lần ở $T_{\text{mem}} = 20$. Gradient tại những
bước gần đầu ra mạnh hơn gradient tại những bước xa đầu ra 25 lần. Đây không
phải là suy luận toán học; đây là con số đo được từ một RNN thật.

\textbf{Thứ hai: loss không giảm.} Ở $T_{\text{mem}} = 20$, loss sau huấn
luyện (8.37) cao hơn loss trước huấn luyện (7.91). Mạng không những không
học được, mà còn tệ hơn cả trước khi học. Gradient yếu đến mức các bước cập
nhật trọng số thực chất là nhiễu, và nhiễu đó phá hỏng cả những gì mạng đã
có từ khởi tạo ngẫu nhiên.

\textbf{Thứ ba: đây không phải là vấn đề năng lực.} Một RNN với $d_h = 32$
có sức chứa quá đủ để giải bài toán sao chép $T_{\text{mem}} = 20$ (đầu vào
21 bước, mỗi bước 9 chiều). Vấn đề không phải là mạng quá nhỏ; vấn đề là
gradient không đến được nơi cần đến.

\index{copy task|)}%
